\documentclass[12pt,a4paper]{article}
\usepackage[utf8]{inputenc}
\usepackage[french]{babel}
\usepackage[T1]{fontenc}
\usepackage{amsmath,amssymb}
\usepackage{graphicx}
\usepackage{xcolor}
\usepackage{hyperref}
\usepackage{geometry}
\usepackage{tcolorbox}
\usepackage{tikz}
\usepackage{array}
\geometry{margin=2.5cm}

% Définition des couleurs
\definecolor{titrebleu}{RGB}{0,102,204}
\definecolor{defvert}{RGB}{0,153,76}
\definecolor{exempleorange}{RGB}{255,102,0}
\definecolor{algoviolet}{RGB}{153,0,153}
\definecolor{importantrouge}{RGB}{204,0,0}
\definecolor{fondjaune}{RGB}{255,250,205}
\definecolor{fondvert}{RGB}{230,255,230}
\definecolor{fondbleu}{RGB}{230,240,255}

% Boîtes colorées
\newtcolorbox{exercice}{
    colback=fondjaune,
    colframe=exempleorange,
    fonttitle=\bfseries,
    title=Exercice,
    arc=3mm
}

\newtcolorbox{solution}{
    colback=fondvert,
    colframe=defvert,
    fonttitle=\bfseries,
    title=Solution,
    arc=3mm
}

\newtcolorbox{methode}{
    colback=fondbleu,
    colframe=algoviolet,
    fonttitle=\bfseries,
    title=Methode,
    arc=3mm
}

\newtcolorbox{remarque}{
    colback=white,
    colframe=importantrouge,
    fonttitle=\bfseries,
    title=Remarque,
    arc=3mm
}

% Style des sections
\usepackage{titlesec}
\titleformat{\section}
{\color{titrebleu}\normalfont\LARGE\bfseries}
{\color{titrebleu}\thesection}{1em}{}

\titleformat{\subsection}
{\color{defvert}\normalfont\Large\bfseries}
{\color{defvert}\thesubsection}{1em}{}

\title{\textbf{\Huge\color{titrebleu}Corrections des Exercices}\\[0.5em]
\large\color{defvert}Architecture des Systèmes\\
Chapitre 1 : Fonctionnement d'un ordinateur}
\author{\large Alexandre Guitton}
\date{\large Version 2024-2025}

\begin{document}

\maketitle
\tableofcontents
\newpage

\section{Représentation des Valeurs - Entiers Positifs}

\subsection{Exercice 1 : Conversions vers la base 10}

\begin{exercice}
\textbf{Question 1 :} Combien vaut $100_8$ en base 10 ?

\textbf{Question 2 :} Combien vaut $3B4_{12}$ en base 10 ?
\end{exercice}

\begin{methode}
\textbf{Formule générale :}

Pour convertir un nombre $(a_{n-1}a_{n-2}...a_1a_0)_b$ de la base $b$ vers la base 10 :
$$\text{Valeur} = \sum_{i=0}^{n-1} a_i \times b^i$$
\end{methode}

\begin{solution}
\textbf{Question 1 : $100_8$ en base 10}

$$100_8 = 1 \times 8^2 + 0 \times 8^1 + 0 \times 8^0$$
$$= 1 \times 64 + 0 + 0 = 64$$

\textcolor{importantrouge}{\textbf{Réponse : }} $100_8 = 64_{10}$
\end{solution}

\begin{solution}
\textbf{Question 2 : $3B4_{12}$ en base 10}

En base 12, les chiffres vont de 0 à 11, où A = 10 et B = 11.

$$3B4_{12} = 3 \times 12^2 + 11 \times 12^1 + 4 \times 12^0$$
$$= 3 \times 144 + 11 \times 12 + 4 \times 1$$
$$= 432 + 132 + 4 = 568$$

\textcolor{importantrouge}{\textbf{Réponse : }} $3B4_{12} = 568_{10}$
\end{solution}

\newpage
\subsection{Exercice 2 : Conversion en base 2}

\begin{exercice}
\textbf{Question :} Quelle est la valeur en base 2 du nombre 1792 ?
\end{exercice}

\begin{methode}
\textbf{Méthode par divisions successives :}

\begin{enumerate}
    \item Diviser le nombre par 2
    \item Le reste est le chiffre binaire (de droite à gauche)
    \item Le quotient devient le nouveau nombre
    \item Répéter jusqu'à obtenir un quotient de 0
\end{enumerate}

\textbf{Méthode par puissances de 2 :}
\begin{enumerate}
    \item Écrire les puissances de 2
    \item Trouver la plus grande puissance $\leq$ au nombre
    \item Soustraire et recommencer
\end{enumerate}
\end{methode}

\begin{solution}
\textbf{Méthode par divisions successives :}

\begin{center}
\begin{tabular}{|c|c|c|}
\hline
\textbf{Division} & \textbf{Quotient} & \textbf{Reste} \\
\hline
1792 ÷ 2 & 896 & 0 \\
896 ÷ 2 & 448 & 0 \\
448 ÷ 2 & 224 & 0 \\
224 ÷ 2 & 112 & 0 \\
112 ÷ 2 & 56 & 0 \\
56 ÷ 2 & 28 & 0 \\
28 ÷ 2 & 14 & 0 \\
14 ÷ 2 & 7 & 0 \\
7 ÷ 2 & 3 & 1 \\
3 ÷ 2 & 1 & 1 \\
1 ÷ 2 & 0 & 1 \\
\hline
\end{tabular}
\end{center}

En lisant les restes de bas en haut : $11100000000_2$

\textbf{Verification :} $1792 = 1024 + 512 + 256 = 2^{10} + 2^{9} + 2^{8}$

$1792_{10} = \textcolor{importantrouge}{\textbf{11100000000}_2}$
\end{solution}

\newpage
\subsection{Exercice 3 : Conversion en base 5}

\begin{exercice}
\textbf{Question :} Quelle est la valeur en base 5 du nombre 59 ?
\end{exercice}

\begin{solution}
\textbf{Méthode par divisions successives :}

\begin{center}
\begin{tabular}{|c|c|c|}
\hline
\textbf{Division} & \textbf{Quotient} & \textbf{Reste} \\
\hline
59 ÷ 5 & 11 & 4 \\
11 ÷ 5 & 2 & 1 \\
2 ÷ 5 & 0 & 2 \\
\hline
\end{tabular}
\end{center}

En lisant les restes de bas en haut : $214_5$

\textbf{Verification :}
$214_5 = 2 \times 5^2 + 1 \times 5^1 + 4 \times 5^0 = 50 + 5 + 4 = 59$

\textcolor{importantrouge}{\textbf{Reponse : }} $59_{10} = 214_5$
\end{solution}

\newpage
\subsection{Exercice 4 : Conversions binaire - hexadécimal}

\begin{exercice}
\textbf{Question :} Écrire $AF4_{16}$ en base 2, ensuite en base 4.
\end{exercice}

\begin{methode}
\textbf{Conversion hexadécimal → binaire :} Chaque chiffre hexadécimal = 4 bits (car $16 = 2^4$)

\textbf{Conversion binaire → base 4 :} Regrouper par blocs de 2 bits (car $4 = 2^2$)
\end{methode}

\begin{solution}
\textbf{Étape 1 : Conversion hexadécimal → binaire}

\begin{itemize}
    \item A$_{16}$ = 10$_{10}$ = 1010$_2$
    \item F$_{16}$ = 15$_{10}$ = 1111$_2$
    \item 4$_{16}$ = 4$_{10}$ = 0100$_2$
\end{itemize}

$$AF4_{16} = \textcolor{defvert}{\textbf{101011110100}_2}$$

\textbf{Étape 2 : Conversion binaire → base 4}

Regrouper par blocs de 2 bits : $10\ 10\ 11\ 11\ 01\ 00_2$

$$AF4_{16} = \textcolor{algoviolet}{\textbf{223310}_4}$$

\textbf{Verification :} $AF4_{16} = 2560 + 240 + 4 = 2804_{10}$ et $223310_4 = 2804_{10}$
\end{solution}

\newpage
\subsection{Exercice 4 bis : Binaire vers hexadécimal}

\begin{exercice}
\textbf{Question :} Convertir $11010110101_2$ en hexadécimal.
\end{exercice}

\begin{solution}
\textbf{Regrouper par blocs de 4 bits (de droite à gauche) :}

$$11010110101_2 = 0110\ 1011\ 0101_2$$

\textbf{Convertir chaque bloc :}
\begin{itemize}
    \item $0110_2 = 6_{16}$
    \item $1011_2 = B_{16}$
    \item $0101_2 = 5_{16}$
\end{itemize}

$$11010110101_2 = \textcolor{importantrouge}{\textbf{6B5}_{16}}$$
\end{solution}

\newpage
\section{Additions Binaires}

\subsection{Exercice 5 : Additions en base 2}

\begin{exercice}
\textbf{Question 1 :} Calculer $100111_2 + 101110_2$

\textbf{Question 2 :} Calculer $010101_2 + 111100_2$
\end{exercice}

\begin{methode}
\textbf{Addition binaire :}
\begin{itemize}
    \item 0 + 0 = 0
    \item 0 + 1 = 1
    \item 1 + 0 = 1
    \item 1 + 1 = 10 (0 avec retenue de 1)
    \item 1 + 1 + 1 = 11 (1 avec retenue de 1)
\end{itemize}
\end{methode}

\begin{solution}
\textbf{Question 1 : $100111_2 + 101110_2$}

\begin{center}
\begin{tabular}{cccccccc}
  &   & 1 & 0 & 0 & 1 & 1 & 1 \\
+ &   & 1 & 0 & 1 & 1 & 1 & 0 \\
\hline
\textbf{res} & \textbf{1} & \textbf{0} & \textbf{1} & \textbf{0} & \textbf{1} & \textbf{0} & \textbf{1} \\
\end{tabular}
\end{center}

$$100111_2 + 101110_2 = \textcolor{importantrouge}{\textbf{1010101}_2}$$

\textbf{Verification :} $39 + 46 = 85$
\end{solution}

\begin{solution}
\textbf{Question 2 : $010101_2 + 111100_2$}

\begin{center}
\begin{tabular}{cccccccc}
  &   & 0 & 1 & 0 & 1 & 0 & 1 \\
+ &   & 1 & 1 & 1 & 1 & 0 & 0 \\
\hline
\textbf{res} & \textbf{1} & \textbf{0} & \textbf{1} & \textbf{0} & \textbf{0} & \textbf{0} & \textbf{1} \\
\end{tabular}
\end{center}

$$010101_2 + 111100_2 = \textcolor{importantrouge}{\textbf{1010001}_2}$$

\textbf{Verification :} $21 + 60 = 81$
\end{solution}

\newpage
\section{Entiers Naturels - Complément Logique}

\subsection{Exercice 6 : Valeurs en complément logique}

\begin{exercice}
\textbf{Question :} Quelle est la valeur décimale de :
\begin{enumerate}
    \item $00001101_2$ en complément logique
    \item $11101001_2$ en complément logique
\end{enumerate}
\end{exercice}

\begin{methode}
\textbf{Complément logique :}
\begin{itemize}
    \item Si bit de poids fort = 0 → nombre positif, conversion directe
    \item Si bit de poids fort = 1 → nombre négatif, inverser tous les bits puis convertir
\end{itemize}
\end{methode}

\begin{solution}
\textbf{Question 1 : $00001101_2$}

Bit de poids fort = 0 → nombre positif

$$00001101_2 = 8 + 4 + 1 = \textcolor{importantrouge}{\textbf{+13}}$$
\end{solution}

\begin{solution}
\textbf{Question 2 : $11101001_2$}

Bit de poids fort = 1 → nombre négatif

Inversion : $11101001 \rightarrow 00010110_2 = 22_{10}$

$$11101001_2 = \textcolor{importantrouge}{\textbf{-22}}$$
\end{solution}

\newpage
\subsection{Exercice 8 : Calculs en complément logique}

\begin{exercice}
\textbf{Question 1 :} Calculer $7 - 3$ en base 2 sur 4 bits

\textbf{Question 2 :} Calculer $22 - 22$ en base 2 sur 8 bits
\end{exercice}

\begin{methode}
\textbf{Addition en complément logique :}
\begin{enumerate}
    \item Faire l'addition binaire normale
    \item Si une retenue déborde, la réajouter au résultat (retenue circulaire)
\end{enumerate}
\end{methode}

\begin{solution}
\textbf{Question 1 : $7 - 3 = 7 + (-3)$}

Représentations : $7 = 0111_2$, $-3 = 1100_2$ (inversion de $0011_2$)

$$0111 + 1100 = 10011$$

Retenue débordante réinjectée : $0011 + 1 = 0100_2 = 4$

$$7 - 3 = \textcolor{importantrouge}{\textbf{4}}$$
\end{solution}

\begin{solution}
\textbf{Question 2 : $22 - 22$}

$22 = 00010110_2$, $-22 = 11101001_2$

Addition donne $11111111_2$ + retenue débordante

$$11111111 + 1 = 100000000 \rightarrow 00000000 = -0$$

\textbf{Problème :} En complément logique, 0 et -0 sont différents!
\end{solution}

\newpage
\section{Entiers Naturels - Complément Arithmétique}

\subsection{Exercice 7 : Écriture en complément arithmétique}

\begin{exercice}
\textbf{Question 1 :} Écrire -17 en complément arithmétique sur 8 bits

\textbf{Question 2 :} Écrire -0 en complément arithmétique sur 4 bits
\end{exercice}

\begin{methode}
\textbf{Complément arithmétique (complément à 2) :}
\begin{enumerate}
    \item Écrire $n$ en binaire
    \item Inverser tous les bits
    \item Ajouter 1
\end{enumerate}
\end{methode}

\begin{solution}
\textbf{Question 1 : -17 sur 8 bits}

$$17 = 00010001 \rightarrow 11101110 \rightarrow 11101111$$

$$-17 = \textcolor{importantrouge}{\textbf{11101111}_2}$$
\end{solution}

\begin{solution}
\textbf{Question 2 : -0 sur 4 bits}

$$0 = 0000 \rightarrow 1111 \rightarrow 10000 \rightarrow 0000$$

$$-0 = \textcolor{importantrouge}{\textbf{0000}_2}$$

\textbf{Avantage :} Une seule représentation pour 0!
\end{solution}

\newpage
\subsection{Exercice 9 : Calculs en complément arithmétique}

\begin{exercice}
\textbf{Question 1 :} Calculer $6 - 5$ sur 4 bits

\textbf{Question 2 :} Calculer $3 - 7$ sur 4 bits
\end{exercice}

\begin{methode}
\textbf{Addition en complément arithmétique :}
\begin{enumerate}
    \item Faire l'addition binaire
    \item Ignorer la retenue débordante
\end{enumerate}
\end{methode}

\begin{solution}
\textbf{Question 1 : $6 - 5 = 6 + (-5)$}

$6 = 0110$, $-5 = 1011$ (CA de 5)

$$0110 + 1011 = 10001 \rightarrow 0001 = +1$$
\end{solution}

\begin{solution}
\textbf{Question 2 : $3 - 7 = 3 + (-7)$}

$3 = 0011$, $-7 = 1001$ (CA de 7)

$$0011 + 1001 = 1100$$

Bit poids fort = 1 → négatif. Valeur : $-4$

$$3 - 7 = \textcolor{importantrouge}{\textbf{-4}}$$
\end{solution}

\newpage
\subsection{Exercice 9 bis : Débordements}

\begin{exercice}
\textbf{Question :} Calculer $5 + 6$ et $-5 + (-6)$ sur 4 bits. Y a-t-il débordement ?
\end{exercice}

\begin{remarque}
Sur 4 bits en CA : valeurs de $-8$ à $+7$

Débordement si : positif + positif = négatif OU négatif + négatif = positif
\end{remarque}

\begin{solution}
\textbf{$5 + 6$ :} $0101 + 0110 = 1011 = -5$

\textcolor{importantrouge}{\textbf{Débordement !}} Résultat $11 > 7$
\end{solution}

\begin{solution}
\textbf{$-5 + (-6)$ :} $1011 + 1010 = 10101 \rightarrow 0101 = +5$

\textcolor{importantrouge}{\textbf{Débordement !}} Résultat $-11 < -8$
\end{solution}

\newpage
\section{Nombres Réels - IEEE 754}

\subsection{Exercice 10 : Encodage IEEE 754}

\begin{exercice}
\textbf{Question :} Encoder -0,1015625 en simple précision IEEE 754.
\end{exercice}

\begin{methode}
\textbf{Format IEEE 754 simple précision :}
\begin{itemize}
    \item 1 bit de signe
    \item 8 bits d'exposant (biaisé de 127)
    \item 23 bits de mantisse (bit caché)
\end{itemize}
\end{methode}

\begin{solution}
\textbf{Signe :} 1 (négatif)

\textbf{Conversion :} $0,1015625 = 0,0001101_2 = 1,101 \times 2^{-4}$

\textbf{Exposant :} $-4 + 127 = 123 = 01111011_2$

\textbf{Mantisse :} 101 sur 23 bits = 10100000000000000000000

$$-0,1015625 = \textcolor{importantrouge}{\textbf{1\ 01111011\ 10100000000000000000000}}$$
\end{solution}

\newpage
\section{Circuits Logiques}

\subsection{Exercice 12 : Simplification de fonction logique}

\begin{exercice}
\textbf{Question :} Simplifier $L = \bar{S} \cdot \bar{A} \cdot B + \bar{S} \cdot A \cdot \bar{B} + \bar{S} \cdot A \cdot B$

Puis réaliser le circuit simplifié.
\end{exercice}

\begin{methode}
\textbf{Méthodes de simplification :}
\begin{enumerate}
    \item Factorisation
    \item Propriétés de l'algèbre de Boole
    \item Tableaux de Karnaugh
\end{enumerate}

\textbf{Propriétés utiles :}
\begin{itemize}
    \item $A + A = A$ (idempotence)
    \item $A \cdot \bar{A} = 0$
    \item $A + \bar{A} = 1$
    \item $A \cdot (B + C) = A \cdot B + A \cdot C$ (distributivité)
\end{itemize}
\end{methode}

\begin{solution}
\textbf{Simplification par factorisation :}

$L = \bar{S} \cdot \bar{A} \cdot B + \bar{S} \cdot A \cdot \bar{B} + \bar{S} \cdot A \cdot B$

\textbf{Étape 1 :} Factoriser par $\bar{S}$

$L = \bar{S} \cdot (\bar{A} \cdot B + A \cdot \bar{B} + A \cdot B)$

\textbf{Étape 2 :} Factoriser les deux derniers termes par $A$

$L = \bar{S} \cdot (\bar{A} \cdot B + A \cdot (\bar{B} + B))$

\textbf{Étape 3 :} Simplifier $\bar{B} + B = 1$

$L = \bar{S} \cdot (\bar{A} \cdot B + A \cdot 1)$
$L = \bar{S} \cdot (\bar{A} \cdot B + A)$

\textbf{Étape 4 :} Utiliser $A + \bar{A} \cdot B = A + B$ (théorème d'absorption)

$\boxed{L = \textcolor{importantrouge}{\bar{S} \cdot (A + B)}}$

\textbf{Comparaison :}
\begin{itemize}
    \item Circuit original : 3 portes NOT, 3 portes AND à 3 entrées, 1 porte OR
    \item Circuit simplifié : 1 porte NOT, 1 porte OR, 1 porte AND
    \item \textcolor{defvert}{\textbf{Gain énorme !}}
\end{itemize}
\end{solution}

\newpage
\subsection{Exercice 13 : Simplification par Karnaugh (3 variables)}

\begin{exercice}
\textbf{Question :} Simplifier la fonction suivante à l'aide d'un tableau de Karnaugh :
$F(A,B,C) = \bar{A}\bar{B}C + \bar{A}BC + A\bar{B}C + ABC$
\end{exercice}

\begin{methode}
\textbf{Méthode des tableaux de Karnaugh :}

\begin{enumerate}
    \item Construire le tableau (2 variables en ligne, 1 en colonne)
    \item Remplir avec les valeurs de la fonction
    \item Regrouper les 1 adjacents par blocs de $2^n$ (1, 2, 4, 8...)
    \item Pour chaque regroupement, identifier les variables constantes
\end{enumerate}

\textbf{Règle :} Plus les regroupements sont grands, plus la simplification est efficace.
\end{methode}

\begin{solution}
\textbf{Étape 1 : Table de vérité}

\begin{center}
\begin{tabular}{|c|c|c|c|}
\hline
\textbf{A} & \textbf{B} & \textbf{C} & \textbf{F} \\
\hline
0 & 0 & 0 & 0 \\
0 & 0 & 1 & 1 \\
0 & 1 & 0 & 0 \\
0 & 1 & 1 & 1 \\
1 & 0 & 0 & 0 \\
1 & 0 & 1 & 1 \\
1 & 1 & 0 & 0 \\
1 & 1 & 1 & 1 \\
\hline
\end{tabular}
\end{center}

\textbf{Étape 2 : Tableau de Karnaugh}

\begin{center}
\begin{tabular}{|c|c|c|}
\hline
\textbf{AB/C} & \textbf{0} & \textbf{1} \\
\hline
00 & 0 & 1 \\
\hline
01 & 0 & 1 \\
\hline
11 & 0 & 1 \\
\hline
10 & 0 & 1 \\
\hline
\end{tabular}
\end{center}

\textbf{Étape 3 : Regroupement}

Toute la colonne $C=1$ peut être regroupée → Variable constante : $C$

\textbf{Fonction simplifiée :}

$F(A,B,C) = \textcolor{importantrouge}{\textbf{C}}$

\textbf{Vérification :} $F = 1$ si et seulement si $C = 1$
\end{solution}

\newpage
\subsection{Exercice 14 : Karnaugh à 4 variables}

\begin{exercice}
\textbf{Question :} Simplifier la fonction à 4 variables :
$G(A,B,C,D) = \bar{A}\bar{B}\bar{C}D + \bar{A}\bar{B}CD + \bar{A}B\bar{C}D + \bar{A}BCD + A\bar{B}\bar{C}D + A\bar{B}CD$
\end{exercice}

\begin{solution}
\textbf{Tableau de Karnaugh 4 variables :}

\begin{center}
\begin{tabular}{|c|c|c|c|c|}
\hline
\textbf{AB/CD} & \textbf{00} & \textbf{01} & \textbf{11} & \textbf{10} \\
\hline
00 & 0 & 1 & 1 & 0 \\
\hline
01 & 0 & 1 & 1 & 0 \\
\hline
11 & 0 & 0 & 0 & 0 \\
\hline
10 & 0 & 1 & 1 & 0 \\
\hline
\end{tabular}
\end{center}

\textbf{Regroupements :}

\textbf{Regroupement 1 :} Rectangle $2 \times 2$ en haut (lignes 00, 01)
\begin{itemize}
    \item Variables constantes : $\bar{A}$ et $D$
    \item Terme : $\bar{A} \cdot D$
\end{itemize}

\textbf{Regroupement 2 :} Rectangle $1 \times 2$ en bas (ligne 10)
\begin{itemize}
    \item Variables constantes : $A$, $\bar{B}$ et $D$
    \item Terme : $A \cdot \bar{B} \cdot D$
\end{itemize}

\textbf{Fonction simplifiée :}

$G(A,B,C,D) = \bar{A}D + A\bar{B}D$

Ou factorisé par $D$ :

$G(A,B,C,D) = \textcolor{importantrouge}{D(\bar{A} + A\bar{B})}$

En utilisant $\bar{A} + A\bar{B} = \bar{A} + \bar{B}$ :

$G(A,B,C,D) = \textcolor{importantrouge}{D(\bar{A} + \bar{B})}$
\end{solution}

\newpage
\subsection{Exercice 15 : Conception d'un demi-additionneur}

\begin{exercice}
\textbf{Question :} Concevoir un circuit qui additionne deux bits $A$ et $B$, produisant une somme $S$ et une retenue $C$ (carry).
\end{exercice}

\begin{solution}
\textbf{Étape 1 : Table de vérité}

\begin{center}
\begin{tabular}{|c|c|c|c|}
\hline
\textbf{A} & \textbf{B} & \textbf{S (Somme)} & \textbf{C (Retenue)} \\
\hline
0 & 0 & 0 & 0 \\
0 & 1 & 1 & 0 \\
1 & 0 & 1 & 0 \\
1 & 1 & 0 & 1 \\
\hline
\end{tabular}
\end{center}

\textbf{Étape 2 : Fonctions logiques}

Pour la somme $S$ :

$S = \bar{A}B + A\bar{B} = \textcolor{importantrouge}{A \oplus B}$ \textbf{(XOR)}

Pour la retenue $C$ :

$C = \textcolor{importantrouge}{A \cdot B}$ \textbf{(AND)}

\textbf{Étape 3 : Description du circuit}

Le circuit nécessite :
\begin{itemize}
    \item 1 porte XOR pour la somme $S$
    \item 1 porte AND pour la retenue $C$
\end{itemize}

\textbf{Nom :} Ce circuit s'appelle un \textbf{demi-additionneur} (half-adder).

\textbf{Remarque :} Pour un additionneur complet (full-adder), il faudrait aussi prendre en compte une retenue entrante.
\end{solution}

\newpage
\section{Résumé et Conseils}

\begin{remarque}
\textbf{Points clés à retenir :}

\textbf{1. Conversions de bases}
\begin{itemize}
    \item Base $b$ vers base 10 : $\sum a_i \times b^i$
    \item Base 10 vers base $b$ : divisions successives
    \item Base 2 ↔ Base 16 : grouper par 4 bits
    \item Base 2 ↔ Base 4 : grouper par 2 bits
\end{itemize}

\textbf{2. Additions binaires}
\begin{itemize}
    \item 0 + 0 = 0
    \item 0 + 1 = 1
    \item 1 + 1 = 10 (retenue)
    \item 1 + 1 + 1 = 11 (retenue)
\end{itemize}

\textbf{3. Complément logique}
\begin{itemize}
    \item Pour $-n$ : inverser tous les bits de $n$
    \item Addition : retenue débordante réinjectée
    \item Problème : deux représentations pour 0
\end{itemize}

\textbf{4. Complément arithmétique}
\begin{itemize}
    \item Pour $-n$ : inverser les bits de $n$ puis ajouter 1
    \item Addition : ignorer la retenue finale
    \item Avantage : une seule représentation pour 0
    \item Sur $n$ bits : de $-2^{n-1}$ à $2^{n-1}-1$
    \item Débordement : positif + positif = négatif OU négatif + négatif = positif
\end{itemize}

\textbf{5. IEEE 754}
\begin{itemize}
    \item Simple précision : 1 bit signe - 8 bits exposant - 23 bits mantisse
    \item Exposant biaisé de 127
    \item Mantisse normalisée : $1.m \times 2^e$ (bit caché)
    \item Valeurs spéciales : zéro (exp=0), infini (exp=255, mant=0), NaN (exp=255, mant$\neq$0)
\end{itemize}

\textbf{6. Caractères}
\begin{itemize}
    \item ASCII : 7 bits (128 caractères)
    \item 'A' = 65, 'a' = 97 (différence de 32)
    \item Chiffres : '0' = 48 à '9' = 57
\end{itemize}

\textbf{7. Circuits logiques}
\begin{itemize}
    \item Portes de base : NOT, AND, OR
    \item Autres portes : XOR, NAND, NOR
    \item Simplification : factorisation et propriétés booléennes
    \item Tableaux de Karnaugh : regrouper les 1 adjacents
    \item Théorème d'absorption : $A + \bar{A} \cdot B = A + B$
\end{itemize}
\end{remarque}

\newpage
\begin{methode}
\textbf{Checklist pour vérifier vos réponses :}

\begin{enumerate}
    \item \textbf{Conversions :} toujours vérifier en convertissant dans l'autre sens
    \item \textbf{Additions binaires :} vérifier en base 10
    \item \textbf{Complément logique :} le bit de poids fort indique le signe
    \item \textbf{Complément arithmétique :} vérifier les débordements
    \item \textbf{IEEE 754 :} vérifier l'exposant biaisé (127 pour simple précision)
    \item \textbf{IEEE 754 :} ne pas oublier le bit caché de la mantisse
    \item \textbf{ASCII :} vérifier que les codes sont dans la plage 0-127
    \item \textbf{Circuits :} vérifier avec la table de vérité
    \item \textbf{Karnaugh :} vérifier que les regroupements sont des puissances de 2
\end{enumerate}
\end{methode}

\begin{remarque}
\textbf{Erreurs fréquentes à éviter :}

\begin{itemize}
    \item \textbf{Conversions :} Oublier de lire les restes de bas en haut
    \item \textbf{Hexadécimal :} Confondre les lettres (A=10, B=11, etc.)
    \item \textbf{Complément arithmétique :} Oublier d'ajouter 1 après l'inversion
    \item \textbf{IEEE 754 :} Oublier le biais de 127 pour l'exposant
    \item \textbf{IEEE 754 :} Oublier le bit caché (le 1 initial de la mantisse)
    \item \textbf{Karnaugh :} Mal organiser le tableau (utiliser l'ordre Gray : 00, 01, 11, 10)
    \item \textbf{Débordement :} Ne pas vérifier si le résultat est représentable
\end{itemize}
\end{remarque}

\begin{methode}
\textbf{Astuces pour gagner du temps :}

\begin{itemize}
    \item \textbf{Binaire ↔ Hexadécimal :} Mémoriser la table de conversion (0000=0 à 1111=F)
    \item \textbf{Puissances de 2 :} Mémoriser jusqu'à $2^{10}=1024$
    \item \textbf{Complément arithmétique :} Pour trouver $-n$, partir de droite, recopier jusqu'au premier 1 inclus, puis inverser le reste
    \item \textbf{IEEE 754 :} Pour les nombres simples, penser aux puissances de 2
    \item \textbf{Karnaugh :} Chercher d'abord les plus grands regroupements
\end{itemize}
\end{methode}

\newpage
\section{Annexes}

\subsection*{Table des puissances de 2}

\begin{center}
\begin{tabular}{|c|c||c|c|}
\hline
\textbf{Exposant} & \textbf{Valeur} & \textbf{Exposant} & \textbf{Valeur} \\
\hline
$2^0$ & 1 & $2^{10}$ & 1024 \\
$2^1$ & 2 & $2^{11}$ & 2048 \\
$2^2$ & 4 & $2^{12}$ & 4096 \\
$2^3$ & 8 & $2^{13}$ & 8192 \\
$2^4$ & 16 & $2^{14}$ & 16384 \\
$2^5$ & 32 & $2^{15}$ & 32768 \\
$2^6$ & 64 & $2^{16}$ & 65536 \\
$2^7$ & 128 & $2^{20}$ & 1048576 \\
$2^8$ & 256 & $2^{30}$ & 1073741824 \\
$2^9$ & 512 & $2^{31}$ & 2147483648 \\
\hline
\end{tabular}
\end{center}

\subsection*{Table de conversion Binaire - Hexadécimal}

\begin{center}
\begin{tabular}{|c|c|c||c|c|c|}
\hline
\textbf{Déc.} & \textbf{Bin.} & \textbf{Hex.} & \textbf{Déc.} & \textbf{Bin.} & \textbf{Hex.} \\
\hline
0 & 0000 & 0 & 8 & 1000 & 8 \\
1 & 0001 & 1 & 9 & 1001 & 9 \\
2 & 0010 & 2 & 10 & 1010 & A \\
3 & 0011 & 3 & 11 & 1011 & B \\
4 & 0100 & 4 & 12 & 1100 & C \\
5 & 0101 & 5 & 13 & 1101 & D \\
6 & 0110 & 6 & 14 & 1110 & E \\
7 & 0111 & 7 & 15 & 1111 & F \\
\hline
\end{tabular}
\end{center}

\subsection*{Propriétés de l'algèbre de Boole}

\begin{center}
\begin{tabular}{|l|l|}
\hline
\textbf{Propriété} & \textbf{Formule} \\
\hline
Identité & $A + 0 = A$ \quad $A \cdot 1 = A$ \\
Annulation & $A + 1 = 1$ \quad $A \cdot 0 = 0$ \\
Idempotence & $A + A = A$ \quad $A \cdot A = A$ \\
Complémentation & $A + \bar{A} = 1$ \quad $A \cdot \bar{A} = 0$ \\
Double négation & $\overline{\bar{A}} = A$ \\
Commutativité & $A + B = B + A$ \quad $A \cdot B = B \cdot A$ \\
Associativité & $(A + B) + C = A + (B + C)$ \\
Distributivité & $A \cdot (B + C) = A \cdot B + A \cdot C$ \\
Absorption & $A + A \cdot B = A$ \quad $A \cdot (A + B) = A$ \\
& $A + \bar{A} \cdot B = A + B$ \\
De Morgan & $\overline{A + B} = \bar{A} \cdot \bar{B}$ \\
& $\overline{A \cdot B} = \bar{A} + \bar{B}$ \\
\hline
\end{tabular}
\end{center}

\subsection*{Codes ASCII principaux}

\begin{center}
\begin{tabular}{|c|c||c|c|}
\hline
\textbf{Caractère} & \textbf{Code décimal} & \textbf{Caractère} & \textbf{Code décimal} \\
\hline
'0' & 48 & 'A' & 65 \\
'1' & 49 & 'B' & 66 \\
'2' & 50 & 'C' & 67 \\
... & ... & ... & ... \\
'9' & 57 & 'Z' & 90 \\
\hline
espace & 32 & 'a' & 97 \\
'!' & 33 & 'b' & 98 \\
'+' & 43 & 'c' & 99 \\
'/' & 47 & ... & ... \\
':' & 58 & 'z' & 122 \\
\hline
\end{tabular}
\end{center}

\vspace{1cm}

\begin{center}
\textcolor{titrebleu}{\Large\textbf{Fin du document}}

\vspace{0.5cm}

\textcolor{defvert}{\large Bonne révision et bon courage pour vos examens !}
\end{center}

\end{document}